\documentclass[12pt]{article}

\usepackage[utf8]{inputenc}
\usepackage[T1]{fontenc}
\usepackage[romanian]{babel}
\usepackage{amsmath, amssymb, amsthm}
\usepackage{geometry}
\geometry{a4paper, margin=2.5cm}

\title{\Large\textbf{Propoziția lui Weierstrass privind derivarea sub integrală}}
\date{}

\theoremstyle{plain}
\newtheorem{proposition}{Propoziție}

\begin{document}

\maketitle

\begin{proposition}[Weierstrass]
Fie $\Omega$ o mulțime deschisă din $\mathbb{C}$ și fie $[a,b]$ un interval închis.  
Considerăm o funcție


\[
F : \Omega \times [a,b] \to \mathbb{C},
\]


continuă pe $\Omega \times [a,b]$, iar pentru fiecare $t \in [a,b]$ funcția


\[
z \mapsto F(z,t)
\]


este olomorfă în $\Omega$.

Definim


\[
G(z) = \int_a^b F(z,t)\, dt.
\]



Atunci $G$ este olomorfă în $\Omega$, iar pentru orice $z \in \Omega$ are loc


\[
G'(z) = \int_a^b \frac{\partial F}{\partial z}(z,t)\, dt.
\]


\end{proposition}

\end{document}
